\clearpage

\section{Discussion and Conclusions}\label{sec-nbcp-discussion}

\subsection{Consequence for the paper and the next
calculation}\label{consequence-for-the-paper-and-the-next-calculation}

The threefold argument is appropriate for a generic V branch with
nonzero \(J_\Gamma\) and a nonvanishing threefold coefficient. It does
not apply unchanged on the exactly \(J_\Gamma=0\) axis. There, the
pure-PD V branch has the sixfold selection rule and an intermediate
algebraic phase is symmetry-allowed within the stated phase-only
assumptions. This is a necessary qualification of the domain of the
paper's thermal V-state conclusion; it does not establish that a
microscopic V supersolid actually survives at finite temperature.

The zero-SOC Y calculation fixes a classical normalization baseline, and
the SOC extension identifies internal relaxation and a mixed dynamical
term. The full-zone harmonic screen at 0.2 T finds no instability for
the sampled weak-SOC Y backgrounds, including the saved selected angles
on the individual SOC axes. The subsequent angular matching quantifies
the classical tensor variation and the leading vacuum-potential
harmonics, including mixed SOC. The nonlinear smooth-wave test in
Section~\ref{sec-nbcp-y-nonlinear-gradient} further supports its
classical gradient sector at long wavelength, with finite phase
excursions and internal relaxation. The classical density-wall
calculation in Section~\ref{sec-nbcp-y-density-wall} resolves a finite
energy cost on the tested local branches, with microscopic cores and
metastable phase-offset responses. Classical vortex cores, vortex pairs
and wall junctions are now computed
(Sections~\ref{sec-nbcp-matching-results}
and~\ref{sec-nbcp-y-density-wall}). The open input is the quantum
angular potential itself. At \(S=1/2\) its next order in \(1/S\) is
19--27 times the leading angular curvature, with the opposite sign
(Appendix Section~\ref{sec-nbcp-gap-validity}). The pinning strength,
the selected orientation, the pseudo-Goldstone gaps and the lock-in
scale therefore need a calculation that does not rely on the \(1/S\)
expansion. A classical spin calculation must still be kept separate
from a phase model that uses a quantum potential. The bond-displacement convention of the code is fixed, so
odd-in-momentum responses refer to physical momentum; assigning them
laboratory directions requires only the orientation of the model \(x\)
axis relative to the crystal axes, which is a material input.
Quantum gradient corrections and the V calculation remain separate work
before a general continuum theory can be assembled. Near the PD axis in
V, the threefold coefficient is linear in \(J_\Gamma\)
(Section~\ref{sec-nbcp-v-mixed-soc}). It competes with the sixfold term
and closes the leading-order gap at \(J_\Gamma/J_{\rm PD}\approx0.09\)--0.22
for the sampled \(J_{\rm PD}\). A thermal study near that axis must
therefore retain both harmonics and the unequal barriers they produce.
Finite-size scaling or a coupled-defect analysis is still needed before
assigning a thermal transition or a supersolid phase.

\subsection{Scope of the present
evidence}\label{scope-of-the-present-evidence}

The available results connect a microscopic nearest-neighbor Hamiltonian
to local semiclassical pinning and a tested classical Y gradient sector.
They also resolve density-wall profiles within constrained classical
local minima. They leave open the global competition among phases,
vortex and composite defects, quantum gradient corrections and thermal
matching. Skyrmion and V supersolidity therefore retain the explicit
evidence limits stated in their chapters.

Numerical agreement within the same Hamiltonian is evidence for
implementation consistency, not an independent validation of the
material model. The
\href{../../../GOVERNMENT/Working-Pad/issue-notes/open/260918-nbcp-physics-code-review.md}{physics-to-code
review record} maps the claims to implementation and saved checks,
separates previously verified results from open convention issues, and
records the next independent tests. This document remains in review.
